\documentclass[10pt,a4paper]{amsart}
\usepackage[margin=19mm]{geometry}
\usepackage{amsmath,amssymb,mathtools}
\usepackage[scr=rsfs,cal=euler]{mathalfa}
\usepackage{xfrac}
\usepackage[unicode,hidelinks,bookmarks=false]{hyperref}
\newcommand{\ie}{i.e.\ }
\newcommand{\cf}{cf.\ }
\newcommand{\resp}{respectively\ }
\newcommand{\Subsection}{\subsection}
\providecommand{\nobreakdash}{\nobreak\hbox{-}}
\setlength{\emergencystretch}{2em}
\multlinegap0pt
\usepackage{bm}
\usepackage{bbm}
\usepackage{dsfont}
\usepackage{xspace}
\usepackage{cancel}
\usepackage{mathtools}
\usepackage{amsthm}
\makeatletter
\@ifundefined{theorem}{\newtheorem{theorem}{Theorem}}{}
\@ifundefined{lemma}{\newtheorem{lemma}[theorem]{Lemma}}{}
\makeatother
\makeatletter
\expandafter\ifx\csname rezabib\endcsname\relax
\newenvironment{rezabib} 
{\bibdiv\biblist\setupbib} 
{\endbiblist\endbibdiv} 
\fi
\makeatother
\title[Number Fields With Large P\'olya Groups]{Number Fields With Large P\'olya Groups}
\author{Amir Akbary, Abbas Maarefparvar}
\date{Source: arXiv:2508.11125v1 (CC BY 4.0)}
\begin{document}
\maketitle

\noindent\emph{Source and licence.} Number Fields With Large P\'olya Groups. Version arXiv:2508.11125v1 (CC BY 4.0), distributed under the Creative Commons Attribution 4.0 International licence, as recorded in the arXiv licence metadata for this exact version and shown on the abstract page https://arxiv.org/abs/2508.11125v1. Full rights notice: licenses/arxiv-2508.11125-CC-BY-4.0.txt.\par\smallskip

\noindent\textit{Editorial scope.} Counted unit: the Lemma whose source label is \texttt{lemma, log R\_K less than log dK} in the article (source0.tex lines 1426--1428, source label \texttt{lemma, log R\_K less than log dK}), delivered together with the article's own complete original proof of it (source0.tex lines 1431--1506, one \texttt{proof} environment of 76 lines). No mathematics is written, repaired or re-derived: statement and proof are the article's own text line for line. Required context retained uncounted: lemma, Degert (lines 1342--1356). Every definition, display and label of those lines is retained unchanged. Omitted from this package and not redistributed: 1--1341, 1357--1425, 1429--1430, 1507--2407. Nothing inside a retained excerpt is omitted or altered: every retained body is delivered exactly as the article's lines are written, so no omission receipt of any kind accompanies this package. Citations: the retained excerpts carry no citation command at all, so no bibliography entry of the article is transcribed and no reference is redistributed. Reference adaptation: none was needed and none was made; every reference inside the delivered unit resolves inside it, so no unresolved reference or citation is produced. Printed slips kept literal: no printed word, symbol, hypothesis or number of the counted statement or of its proof was altered. Layout and class: the article's own class is \texttt{amsart} with packages including scrextend, amsmath, amssymb, amsfonts, amsthm, amscd, amstext, amsxtra, amsopn, array, url, mathrsfs, enumerate, anysize; the wrapper is the lane's pinned \texttt{amsart} editorial wrapper at 10pt on A4, which restates the theorem-like environments the retained text needs and adds the article's own macro definitions that the retained text needs (none), with extra wrapper packages (bm, bbm, dsfont, xspace, cancel, mathtools). The delivered numbering is the unit's own. Assets: no retained span contains a figure environment, a graphics-inclusion command, a PDF-page inclusion, a table or a TikZ picture; no image file is copied into this package, the package declares no asset dependency and no image file is redistributed with it. Licence: version arXiv:2508.11125v1 of this article is distributed under the Creative Commons Attribution 4.0 International licence, as recorded in the arXiv licence metadata for this exact version and shown on the abstract page https://arxiv.org/abs/2508.11125v1. A full rights notice is delivered at licenses/arxiv-2508.11125-CC-BY-4.0.txt. Characters: every retained excerpt is pure ASCII, so the delivered engine, XeLaTeX with its default Latin Modern text font, reports no missing character.
\begin{lemma}  \label{lemma, Degert}
If $K=\mathbb{Q}(\sqrt{\ell^2+r})$ is a real quadratic field of R-D type, then
	\begin{equation} \label{equation, f.u of R-D type} 
	u_K :=	\left\{
	\begin{array}{ll}
		\ell + \sqrt{\ell^2+r}, \, \, & \textit{if} ~~ |r|= 1,\\
		& \\
		\frac{	\ell + \sqrt{\ell^2+r}}{2}, \, \, & \textit{if} ~~ |r|= 4,\\
		& \\
		\frac{	2\ell^2+r + 2 \ell \sqrt{\ell^2+r}}{|r|}, \, \,  & \textit{if} ~~ |r|\neq 1,4,\\
	\end{array}
	\right.
\end{equation}
is the fundamental unit of $K$. More generally, $u_K$ is a unit if $K$ is of extended R-D type.
\end{lemma}

\begin{lemma}\label{lemma, log R_K less than log dK}
Let $K=\mathbb{Q}(\sqrt{D})$ be a real quadratic field of extended R-D type. Then $ R_K <\log (3 D)$, where $R_K$ denotes the regulator of $K$. 
 \end{lemma}

 \begin{proof} 	
 As usual, denote by $\epsilon_K$ the fundamental unit of $K$. First observe that for $D=5$, we have $\epsilon_K=\frac{-1+\sqrt{5}}{2}$. So, $	R_K=\log\left( \epsilon_K \right) \approx 0.482 < \log(3 D) \approx 2.7$.
 	%\begin{equation*}
 	%	R_K=\log\left(\frac{-1+\sqrt{5}}{2}\right) \approx 0.482 < \log(3 D)=\log (15) \approx 2.7.
 %	\end{equation*}


Next, let $D=\ell^2+r \neq 5$ with $r \mid 4\ell$.  
We have three cases:

	
\noindent 	
\textit{Case 1.} If $|r|=1$, then by Lemma \ref{lemma, Degert}, 
 		\begin{equation*}
 			\epsilon_K =\ell + \sqrt{\ell^2 \pm 1} < \left(2\ell^2 \pm 2\right)+\left(\ell^2 \pm 1\right)=3D.
 		\end{equation*}
 %Similarly, for $r=-1$,
 %\begin{equation*}
 %	\epsilon_K =\ell + \sqrt{\ell^2-1} < \left(2\ell^2-2\right)+\left(\ell^2-1\right)=3D,
 %\end{equation*}
 (Note that $\ell \geq 2$ for $r=-1$.)
 	 
\noindent
\textit{Case 2.} If $|r|=4$, then, by Lemma \ref{lemma, Degert},
	\begin{equation*}
		u_K=\frac{	\ell + \sqrt{\ell^2+r}}{2}
	\end{equation*}
	is a unit of $K$. 
		Hence, 
	\begin{equation*}
		\epsilon_K \leq 	\frac{	\ell + \sqrt{\ell^2 \pm 4}}{2} 	< \left(2 \ell^2\pm 8\right) + \left(\ell^2\pm 4\right)=3D.
	\end{equation*}
%	Also, for $r=-4$,
%	\begin{equation*}
	%	\epsilon_K \leq 	\frac{	\ell + \sqrt{\ell^2-4}}{2} < \left(2\ell^2-8 \right)+ \left(\ell^2-4\right)=3D,
%	\end{equation*}
	(Note that $\ell \geq 3$ for $r=-4$.)

\noindent 
\textit{Case 3.} if $|r| \neq 1,4$, then by Lemma \ref{lemma, Degert}, 
	\begin{equation*}
		u_K=\frac{	2\ell^2+r + 2 \ell \sqrt{\ell^2+r}}{|r|}
	\end{equation*}
 	is a unit of $K$. If $\ell=1$, the only possibility is to have $r=2$, i.e., $D=3$. In this case, $\epsilon_K=2+\sqrt{3}$ and $	R_K=\log{\epsilon_{K}} \approx 1.31 < \log (3D) \approx 2.2$.
 %	\begin{equation*}
 	%	R_K=\log{(2+\sqrt{3})} \approx 1.31 < \log (3D)=\log(9) \approx 2.2.
 %	\end{equation*}

Now assume that $\ell >1$. Since $r \leq 4 \ell \leq 2 \ell^2$, then $\ell \sqrt{\ell^2+r} <  2 \ell^2$.
 	%	\begin{equation*}
 	%	\ell \sqrt{\ell^2+r} \leq   \ell \sqrt{\ell^2+ 2 \ell^2} < 2 \ell^2.
 %	\end{equation*}
 
 \noindent
For $r>0$,
 		\begin{equation*}
 			\epsilon_K \leq 	\frac{	2\ell^2+r + 2 \ell \sqrt{\ell^2+r}}{|r|} \leq 	\frac{	2\ell^2+r + 2 \ell \sqrt{\ell^2+r}}{2}  < \left(	\ell^2+ \frac{r}{2} \right) + 2 \ell^2 < 3 \ell^2 + 3r=3D.
 		\end{equation*}
 		
  \noindent
 	For $r <0$, we have
 			\begin{equation} \label{equation, epsilon less l^2}
 			\epsilon_K \leq 	\frac{	2\ell^2+r + 2 \ell \sqrt{\ell^2+r}}{|r|} <	\frac{	2\ell^2 + 2 \ell \sqrt{\ell^2}}{2}= 2\ell^2.
 		\end{equation}
 	Since $r \mid 4 \ell$, for $\ell \geq 12$, we get
 			$-3 r \leq 12 \ell \leq \ell^2.$
 		Consequently, from \eqref{equation, epsilon less l^2} we find that for $\ell \geq 12$, the inequality $\epsilon_K < 2 \ell^2 \leq 3 \ell^2+3r=3D$ holds.
 		%\begin{equation*}
 		%	\epsilon_K < 2 \ell^2 \leq 3 \ell^2+3r=3D.
 		%\end{equation*}
 	We numerically check that the relation
 	\begin{equation*}
 		\frac{	2\ell^2+r + 2 \ell \sqrt{\ell^2+r}}{|r|} < 3 \ell^2+3r=3D
 	\end{equation*}
 	also holds for $1 \leq \ell \leq 11$ and $r \mid 4 \ell$ with $r \neq \pm 1, \pm 4$.
 \end{proof}

\end{document}
